\documentclass[11pt]{article}
\usepackage[margin=1in]{geometry}
\usepackage{amsmath,amssymb,amsthm}
\usepackage{booktabs}
\usepackage{array}
\usepackage[colorlinks=true,linkcolor=blue!45!black,citecolor=blue!45!black,urlcolor=blue!45!black]{hyperref}
\usepackage{xcolor}

\hypersetup{pdftitle={Why the forall-PO-Free Fragment Is Decidable},
  pdfauthor={Michael Wessel --- with Claude (Opus 4.8, Anthropic)}}

\newcommand{\ALCIRCC}{\ensuremath{\mathcal{ALCI}_{\mathrm{RCC5}}}}
\newcommand{\DR}{\ensuremath{\mathsf{DR}}}
\newcommand{\PO}{\ensuremath{\mathsf{PO}}}
\newcommand{\PP}{\ensuremath{\mathsf{PP}}}
\newcommand{\PPI}{\ensuremath{\mathsf{PPI}}}
\newcommand{\EQ}{\ensuremath{\mathsf{EQ}}}
\newcommand{\Comp}{\mathrm{comp}}
\newcommand{\cl}{\mathrm{cl}}
\newcommand{\code}[1]{\texttt{#1}}

\setlength{\parindent}{0pt}
\setlength{\parskip}{0.55em}

\title{\textbf{Why the $\forall\PO$-Free Fragment Is Decidable}\\[0.35em]
\large A Guide for the Tableau-Minded}
\author{Michael Wessel --- with Claude (Opus 4.8, Anthropic)}
\date{2026-07-22}

\begin{document}
\maketitle

\begin{center}
\fbox{\parbox{0.92\textwidth}{\itshape
A companion to the $\ALCIRCC$ decidability project, pitched one level
more technical than \code{WHY\_ITS\_HARD.md}: for a reader who knows
basic description logics, tableau calculi, and blocking, but not
automata theory or model theory. The goal is not to walk you through
the proofs --- it is to put enough intuition and enough hard, checkable
facts in front of you that you can honestly say ``yes, I believe
that.''

\smallskip
\textbf{Honesty note (updated 2026-08-06):} the fragment is now
\textbf{certified decidable across three of its four quadrants} --- genuine
\emph{computable} \code{Decidable (Satisfiable C0)} in Lean~4
(\code{formal/POFreeLift.lean}, ${\sim}13{,}600$ lines, \textbf{zero
sorries}) for horizontal ($\exists\DR/\PO/\EQ$), ascending-vertical
($\exists\PP$), and descending-vertical ($\exists\PPI$) concepts, each
non-vacuously witnessed. The fourth quadrant (mixing,
$\exists\PO + \exists\PP$) has its merged certificate, its encoding, and a
complete decision certified on a concrete witness; the \emph{general}
mixed extraction is the one remaining formalization. The keystone is a
\emph{constructive} uniformization (\code{rr\_covers}) --- the vertical
``W2$'$'' is a kernel-checked theorem, not an oracle, because removing
$\forall\PO$ lets the cross-relations coordinate for free. These fragment
theorems are \textbf{unreviewed}, and the \emph{full-logic} decidability
(F6) stays \textbf{open} --- forced by $\forall\PO$, which the fragment
removes. This supersedes the 2026-07-22 ``strongly supported, not
certified'' status for the three general quadrants; see \code{LEAN.md} and
\code{ASSEMBLY\_DESIGN.md} \S\S24--25.}}
\end{center}

\section{The claim}\label{sec:claim}

\textbf{Theorem ($\forall\PO$-free fragment).} Concept satisfiability
for the $\forall\PO$-free fragment of $\ALCIRCC$ --- concepts in NNF
containing \textbf{no subformula of the form $\forall\PO.D$} --- is
decidable, with a computable live-width bound
$K(C_0)=(1+m)^2\,2^{2^m}$, where $m=|\cl(C_0)|$.

Read the fine print, because it is generous. The fragment
\textbf{keeps}:
\begin{itemize}
\item $\exists\PO.D$ --- ``overlaps something satisfying $D$'';
\item $\forall\DR.D$ and $\exists\DR.D$ --- universal \emph{and}
  existential disjointness;
\item all vertical modalities --- $\exists\PP$, $\forall\PP$,
  $\exists\PPI$, $\forall\PPI$, nested to any depth;
\item infinite models --- $\exists\PP.\top\sqcap\forall\PP.\exists\PP.\top$
  is in the fragment and is satisfiable only in models with an infinite
  ascending $\PP$-chain.
\end{itemize}

The \emph{only} thing removed is the universal restriction over partial
overlap. That is a strikingly thin excision, and the whole story of
this document is why the difficulty of the full logic is concentrated
in exactly that one constructor.

A calibration point that matters: the fragment is \textbf{not} a toy
that dodges the hard diagnostics. The concept
\[
C_{\mathrm{force}} \;=\; \exists\PP.(\exists\DR.A)\sqcap\forall\DR.\neg A
\]
is $\forall\PO$-free and \textsc{unsat} --- because
$x\mathbin{\PP}y$ and $y\mathbin{\DR}z$ \emph{force} $x\mathbin{\DR}z$
($\Comp(\PP,\DR)=\{\DR\}$), and then $\forall\DR.\neg A$ kills the
required $z\!:\!A$. Every naive tableau that ignores forced edges
accepts it wrongly. So a decision procedure for the fragment must still
do genuine composition-table reasoning; it just never has to fight the
one battle that has defeated every full-logic attempt.

\section{The setting, and why your tableau instincts don't directly
apply}\label{sec:setting}

Three facts about $\ALCIRCC$ that break the standard playbook.
(Strong-$\EQ$ semantics throughout: $\EQ$ means identity, so between
distinct elements only $\DR,\PO,\PP,\PPI$ occur.)

\paragraph{Fact 1: every model is a complete graph.} The five RCC5
relations are jointly exhaustive and pairwise disjoint. Any two
elements of any model stand in exactly one of them --- there is no
``unrelated.'' In $\mathcal{ALC}$, a tableau builds a \emph{tree} and
elements in different branches simply have nothing to do with each
other. Here, every pair of nodes you ever create carries a relation,
whether you asked for it or not, and the composition table
\[
\rho(x,z)\in\Comp(\rho(x,y),\rho(y,z))\qquad\text{for all }x,y,z
\]
constrains all of them, globally. Adding one node adds $n$ new edges,
each of which must close every triangle it participates in.

\paragraph{Fact 2: no finite model property.}
$\exists\PP.\top\sqcap\forall\PP.\exists\PP.\top$ forces an infinite
strictly ascending chain of regions ($\PP$ is a strict order under
strong $\EQ$ --- a $\PP$-cycle would collapse by composition into
$x\mathbin{\PP}x$, which is impossible). So you cannot hope to
enumerate finite models: satisfiable concepts may have none.

\paragraph{Fact 3: local blocking is unsound or incomplete.} In
$\mathcal{ALC}$, blocking works because ``same label set'' means ``same
future obligations.'' Here a node's future obligations depend on its
relations to \emph{the whole graph built so far} (Fact~1), which no
fixed-arity local summary captures. The project tried subset blocking,
triangle blocking, quadruple types, profile blocking, double blocking
--- each fails on a concrete diagnostic. This is not a technicality; it
is the open problem of the full logic (the width bound F6) in disguise.

So the fragment result has to earn its termination some other way. That
way is the \emph{finite-certificate architecture}, and it is worth
understanding on its own before either proof route, because both routes
are instances of it.

\section{The architecture: finite certificates for infinite
models}\label{sec:architecture}

When models are legitimately infinite, ``decidable'' cannot mean
``search for a model.'' It means:

\begin{quote}
\textbf{Replace ``model'' by ``finite blueprint of a model,'' so that
(a) a computer can enumerate and check blueprints, (b) every blueprint
that checks out really describes a model, and (c) every satisfiable
concept has a blueprint of computably bounded size.}
\end{quote}

Call the blueprint a \textbf{certificate}. A decidability proof in this
architecture always has exactly three parts:
\begin{enumerate}
\item \textbf{A finite certificate format} with finitely many
  \emph{finitely checkable} conditions (type coherence, witness
  presence, triangle closure on the finite object, \dots). Checking a
  candidate certificate is trivially decidable.
\item \textbf{Soundness (unfolding).} If a certificate passes the
  checks, the possibly-infinite structure it \emph{generates} ---
  obtained by unrolling its loops --- is a genuine model of the
  concept. This is the direction that must not cheat: the checks are
  finite, the generated object is infinite, and you must prove the
  finite checks control the infinite unfolding.
\item \textbf{Completeness (extraction $+$ compression).} If the
  concept has \emph{any} model, you can extract from it a certificate
  within the computable size bound. This is where ``run out of colors''
  pigeonhole arguments live --- and where a bound on the certificate's
  size must come from.
\end{enumerate}

Given all three, the decision procedure is embarrassing: enumerate all
certificates up to the bound, check each, answer \textsc{sat} iff one
passes. All the mathematics is in (2) and (3); the algorithm itself is
a for-loop.

Four remarks that make this concrete for a tableau reader.

\paragraph{A blocked tableau \emph{is} a certificate.} When your
$\mathcal{ALC}$ tableau stops because node $y$ is blocked by ancestor
$x$, the finished tree with its blocking pointers is precisely a finite
certificate: the ``model'' it describes is the infinite unraveling in
which $y$'s subtree ``continues like $x$'s.'' Soundness of blocking
$=$ the unfolding lemma. So the architecture is not exotic; it is what
blocking always secretly was.

\paragraph{Loops are generators, not quotients.} One subtlety is
promoted to a design principle here. The blocking loop must be read as
``\emph{the pattern continues like this}'' (unfold the loop into
infinitely many distinct copies), never as ``\emph{$y$ equals $x$}''
(bend the edge back). Under strong $\EQ$ the second reading is
inconsistent for vertical loops --- it would make a region a proper
part of itself. The certificate has a $\PP$-self-loop; the model it
generates has an infinite chain of \emph{distinct} elements. Finite
generator, infinite unfolding, no finite quotient.

\paragraph{Under abstract semantics, model-checking is
triangle-checking.} A gift of the composition-table semantics: an
interpretation is \emph{any} set with a total labelling $\rho$
satisfying (i) $\rho(d,d)=\EQ$ and strong $\EQ$, (ii) converse
coherence, (iii) composition closure. Nothing else. So ``the unfolding
is a model'' reduces to ``every triangle of the unfolding is closed
under the table'' --- a purely combinatorial statement, exactly the
kind of thing one can prove by case analysis on where the three corners
sit, and exactly what is kernel-checked in \code{POFreeLift.lean}. (For
networks that are only partially specified, RCC5's \emph{patchwork
property} --- path-consistent atomic networks are globally consistent,
Renz--Nebel 1999 --- completes the missing edges; for a fully labelled
unfolding the network \emph{is} the model outright.)

\paragraph{Live width: what has to be bounded.} The size bound in (3)
is the whole game. The key observation is that most of a certificate's
edges cost nothing: an edge whose value is \emph{forced} by composition
from other edges (a ``shadow'') need not be stored or chosen --- it is
recoverable. What must be bounded is the number of \textbf{live} edges:
genuinely free choices attached to each part of the certificate. The
conjecture that every satisfiable $\ALCIRCC$ concept admits a
presentation of computably bounded live width is \textbf{F6} --- the
single open keystone of the full problem, machine-formalized as the
hypothesis of a certified conditional theorem. \textbf{The
$\forall\PO$-free fragment is precisely the place where a live width
bound is actually proved} rather than conjectured. That is what ``a
fragment result'' means in this project: F6, inhabited.

With the architecture fixed, both proofs of the fragment theorem can be
stated in one sentence each:
\begin{itemize}
\item \textbf{Route 1 (two-tier quotient):} vertical blocking gives a
  finite set of chain-descriptors; the composition table makes every
  witness's relation to a chain \emph{stabilize}, so one certificate
  edge can honestly summarize infinitely many model edges --- for
  $\DR,\PP,\PPI$ always, and for $\PO$ exactly when no $\forall\PO$
  constraint watches.
\item \textbf{Route 2 (ordered-disjoint normal form):} an RCC5 model
  \emph{is} just a strict order plus a well-behaved disjointness
  relation, with $\PO$ as the unconstrained residual;
  $\forall\PP/\forall\PPI/\forall\DR$ obligations are order- and
  disjointness-closure conditions that propagate deterministically, and
  no propagation ever manufactures a $\forall\PO$ obligation --- so a
  standard tree tableau with blocking terminates.
\end{itemize}

The rest of this document unpacks those two sentences.

\section{The algebraic heart: four forced cells, and the relation that
is never forced}\label{sec:algebra}

Everything in both routes reduces to a handful of cells of the RCC5
composition table. You can check each mentally with blobs, and each is
machine-verified from finite set semantics (and, for those used in the
Lean artifacts, kernel-checked).

\textbf{The four singleton cells among proper relations:}
\[
\begin{aligned}
\Comp(\PP,\PP)&=\{\PP\}   &&\text{nesting is transitive}\\
\Comp(\PPI,\PPI)&=\{\PPI\} &&\text{\dots in both directions}\\
\Comp(\PP,\DR)&=\{\DR\}   &&\text{a part of something disjoint from $e$
                              is disjoint from $e$}\\
\Comp(\DR,\PPI)&=\{\DR\}  &&\text{\dots converse form}
\end{aligned}
\]

The third one deserves a picture: if the table is far from the moon,
the coin \emph{on} the table is far from the moon. Disjointness rains
\textbf{down} into parts, deterministically. Together with transitivity
of nesting, these four cells are the \emph{only} deterministic
machinery the whole algebra offers between distinct proper relations.

\textbf{The two cells that define the problem:}
\[
\Comp(\PP,\PO)=\{\PP,\PO,\DR\},\qquad
\Comp(\PPI,\PO)=\{\PPI,\PO\}\qquad\text{--- not singletons.}
\]

Picture the first: a cake $y$ partially overlaps a plate $e$. Now take
a crumb $x$ of the cake ($x\mathbin{\PP}y$). Where is the crumb
relative to the plate? Maybe on the plate ($\PP$), maybe straddling its
rim ($\PO$), maybe on the tablecloth ($\DR$). \emph{Partial overlap
tells you nothing determinate about parts.} The second cell is the
upward version: if something inside $y$ overlaps $e$, then $y$ overlaps
$e$ or contains it --- two options, not one.

And one global fact you can scan the whole table for:

\paragraph{No cell of the table equals $\{\PO\}$.} Partial overlap is
never the uniquely forced outcome of any single composition. $\DR$ is
forced by two cells, $\PP$ and $\PPI$ by transitivity --- but no
\emph{single} path $x$--$y$--$z$ ever pins $x\mathbin{\PO}z$. $\PO$ is
forced by nothing.

\paragraph{And the dual: $\PO$ is available wherever anything is.} Scan
the table again: every \emph{non}-singleton proper cell
\textbf{contains} $\PO$ --- and this propagates along paths (of the ten
composition-sets reachable by composing along paths of any length,
every non-singleton one contains $\PO$; machine-checked, \code{wp26}).
Moreover every entry of the $\PO$ row and column contains $\PO$, and
$\Comp(\PO,\PO)$ is everything --- so $\PO$, once present, stays
available in every direction, and an all-$\PO$ region constrains
nothing. Put the two facts together and every pair of elements faces a
clean dichotomy: either some composition path \emph{pins} it (a
singleton --- forced, deterministic, a recoverable ``shadow''), or
\textbf{$\PO$ is among its options}. $\PO$ is the algebra's escape
valve: never forced on you, always offered where anything is. The only
way a logic can close that valve is to make $\PO$-edges \emph{illegal}
--- and the only constructor that can do so is $\forall\PO$. That, in
one sentence, is why its removal is worth a decidability theorem: in
the fragment, the escape valve is structurally un-closable, and the
unfolding's unboundedly many unpinned pairs all take the value that is
always permitted and never watched.

\paragraph{Calibration, learned the hard way.} That last fact is about
\emph{single} compositions. The \emph{intersection} of several path
constraints \textsc{can} force a $\PO$ edge --- e.g.
$\Comp(\DR,\PO)\cap\Comp(\PPI,\PO)=\{\DR,\PO,\PP\}\cap\{\PPI,\PO\}=\{\PO\}$.
An earlier version of the project's prose overstated ``$\PO$ is never
forced,'' and a cold review caught it (the 16th, with a concrete
two-path witness). Neither proof route below relies on the overstated
version: both enforce \emph{full} composition closure on every
triangle, which subsumes multi-path forcing. Where ``$\PO$ is loose''
is used below, it is always in the precise single-cell sense above.

\paragraph{Yes, $\PO$ can be forced --- and it is harmless.} Do not
read the caveat as fragility. Even inside the fragment, $\PO$ edges can
be pinned: algebraically by joint paths (above), and logically by the
\emph{allowed} universals clashing out the alternatives. In
\[
C \;=\; \forall\PP.\neg A \sqcap \forall\DR.\neg A \sqcap
        \exists\PP.(\neg A \sqcap \exists\PO.A)
\]
the demanded $e\!:\!A$ sits at $\Comp(\PP,\PO)=\{\PP,\PO,\DR\}$ from
the root; $\forall\PP.\neg A$ strikes $\PP$, $\forall\DR.\neg A$
strikes $\DR$, and $\rho(\mathrm{root},e)=\PO$ is \emph{forced} (and
the concept is satisfiable --- $x=\{1,2\}$, $y=\{1,2,3\}$,
$e=\{2,3,4\}$). So the fragment can steer edges \emph{into} $\PO$. Why
is that safe, when steering was the full logic's downfall? Because a
forced $\PO$ edge is \textbf{doubly inert}. \emph{Logically inert:} the
only constructor that fires on a $\PO$ edge --- at either endpoint;
$\PO$ is self-converse and the safety test is two-sided --- is
$\forall\PO$, and none exists in the input or in any label that
propagation will ever write (\S\ref{sec:route2}). The forced edge
lands, and nothing happens. \emph{Compositionally, it never forces
onward alone:} $\PO$ is an argument of \textbf{no} singleton cell (look
at the four: $\PP;\PP$, $\PP;\DR$, $\PPI;\PPI$, $\DR;\PPI$ --- $\PO$
appears in none), so a $\PO$ edge never single-handedly pins another
edge; it can narrow options in joint configurations, and full triangle
closure tracks exactly that. Contrast forcing $\DR$: a forced $\DR$
edge can fire a waiting $\forall\DR$ at its far endpoint --- that is
literally $C_{\mathrm{force}}$ --- real consequences, but
\emph{deterministic} ones the singleton cells propagate. Forcing $\PO$
has no consequences at all. Steering into the watched relations is
trackable; steering into the unwatched one is a dead end for any
adversarial concept.

Now the punchline of the whole section. A universal restriction
$\forall R.D$ is a \emph{standing obligation}: it must hold of every
$R$-neighbor, including $R$-neighbors that arise \emph{indirectly},
forced by composition, at any distance, at any time. Whether such an
obligation is manageable depends entirely on whether the algebra moves
$R$-neighbors around predictably:
\begin{itemize}
\item \textbf{$\forall\DR.D$ is manageable} because $\DR$ propagates
  deterministically downward ($\Comp(\PP,\DR)=\{\DR\}$): the set of
  $\DR$-neighbors of a region grows in a controlled, forced,
  \emph{monotone} way as you descend the nesting order. You always know
  where the obligation will land next.
\item \textbf{$\forall\PP.D$ and $\forall\PPI.D$ are manageable} for
  the same reason via transitivity.
\item \textbf{$\forall\PO.D$ is the wild one}: $\PO$-neighborhoods are
  not transported determinately in \emph{either} vertical direction
  (the two non-singleton cells), and $\PO$ is never created by any
  single composition --- so the obligation must be re-verified against
  a set of neighbors that the algebra refuses to pin down. The
  recurring failure shape across this project's seventeen-review
  history --- jointly coordinating horizontal edges against universal
  obligations at a gluing step --- is only \emph{hard} for edges no
  singleton cell controls; and $\PO$ is the relation with no singleton
  cell anywhere.
\end{itemize}

The fragment theorem is the statement that once $\forall\PO$ is gone,
the remaining obligations are exactly the ones the four singleton cells
can carry. Two independent routes cash this in.

\section{Route 1: the two-tier quotient (the chain-and-phase
argument)}\label{sec:route1}

{\small(Sources, under \code{papers/}:
\code{two\_tier\_quotient\_ALCIRCC5.tex} --- Claude, April 2026, four
revisions under GPT-5.4 review --- and the fragment specialization
\code{po\_free\_fragment\_ALCIRCC5.tex}. Soundness crux certified
2026-07-18 in \code{formal/POFreeLift.lean}.)}

\subsection{The certificate}\label{ssec:cert}

Since models contain infinite $\PP$-chains, start where the infinity
is. Walk up an infinite chain $d_0\mathbin{\PP}d_1\mathbin{\PP}d_2\cdots$
and record each element's Hintikka type. There are at most $2^m$ types,
so the sequence of types is eventually periodic --- familiar
pigeonhole, exactly the intuition behind blocking. Call the recurring
cyclic pattern of types $(\tau_1,\dots,\tau_p)$ a \textbf{period
descriptor}. There are at most $2^{2^m}$ descriptors.

The certificate (``two-tier quotient'') is:
\begin{itemize}
\item \textbf{Tier 1:} finitely many period descriptors, one
  \textbf{kernel} node per descriptor carrying its cyclic type pattern
  with a $\PP$-self-loop (a generator: unfolds to the infinite chain,
  per \S\ref{sec:architecture});
\item \textbf{Tier 2:} finitely many \textbf{designated witnesses} ---
  at most one per existential demand per descriptor --- with
  \textbf{one constant relation} recorded between each witness and each
  kernel: a single edge ``kernel${}\leftrightarrow e:R$''.
\end{itemize}
Plus finite checks: types are Hintikka, every $\exists$ demand has its
witness, every recorded triangle is composition-closed, every recorded
edge respects the universals in both endpoint types (the Safe test).

\subsection{The crux: when can one edge summarize infinitely
many?}\label{ssec:crux}

Here is the exact point where this differs from $\mathcal{ALC}$
blocking, and where the reader should slow down. When the kernel
unfolds into the infinite chain $d_0\mathbin{\PP}d_1\mathbin{\PP}\cdots$,
the certificate's single edge ``kernel${}\leftrightarrow e:R$'' must
stand for \emph{infinitely many} model edges: $\rho(d_i,e)$ for every
$i$. The certificate records one value; the model needs a value at
every rung. For the summary to be honest, the witness's relation to the
chain must be \textbf{the same at every rung that matters} ---
otherwise a finite certificate with constant interfaces simply cannot
express the model.

Does the algebra cooperate? Watch a fixed external region $e$ against a
growing chain $d_0\subset d_1\subset d_2\subset\cdots$, using
$\Comp(\PPI,\cdot)$ to step up the chain:
\[
\begin{aligned}
\text{from }\DR &:\ \text{next}\in\{\DR,\PO,\PPI\}\\
\text{from }\PO &:\ \text{next}\in\{\PO,\PPI\}\\
\text{from }\PPI&:\ \text{next}\in\{\PPI\}\quad\text{(absorbing)}\\
\text{from }\PP &:\ \text{next}\in\{\PP,\PO,\PPI\}
\end{aligned}
\]
Read it as a picture: the chain elements grow; a fixed $e$ can go from
disjoint, to overlapping, to \emph{swallowed} ($\PPI$) --- and never
backward. The relation \textbf{stabilizes}: after finitely many
monotone switches it is constant forever. This is the
``external-relation stabilization'' lemma, and it is checkable from the
table by eye.

Stabilization alone is not enough --- the certificate needs the
relation to be right at \emph{all} phases, not just eventually. This is
where the singleton cells become load-bearing, one relation at a time:
\begin{itemize}
\item \textbf{$\DR$ --- backward forcing.} Suppose an $\exists\DR$-demand
  is satisfied by witness $e$ at a late chain position $d_i$. Then for
  every earlier $j<i$: $d_j\mathbin{\PP}d_i$ and $d_i\mathbin{\DR}e$
  give $\rho(d_j,e)\in\Comp(\PP,\DR)=\{\DR\}$. \textbf{Forced.} A
  $\DR$-witness picked late enough (past the point where all recurrent
  types have appeared) is automatically a $\DR$-witness at every phase.
  The constant edge ``kernel${}\leftrightarrow e:\DR$'' is honest, for
  free.
\item \textbf{$\PP$ --- same cell, same argument.} If $e$ contains
  $d_i$, then by $\Comp(\PP,\PP)=\{\PP\}$ it contains every $d_j$
  below. Constant edge honest.
\item \textbf{$\PPI$ --- forward absorption.} If $e$ is inside $d_i$,
  it is inside every $d_j$ above ($\Comp(\PPI,\PPI)=\{\PPI\}$).
  Constant edge honest.
\end{itemize}

So for $\DR,\PP,\PPI$ the composition table itself \emph{converts} a
one-position witness into an all-position witness. The finite summary
is not an approximation; it is exact. If you believe those three
singleton cells --- and you can check them with blobs in a margin ---
you believe the constant-interface certificate is sound and complete
for these three relations. That is the ``yes!'' moment for Route~1.

\subsection{Where $\PO$ breaks it --- a concrete satisfiable
villain}\label{ssec:villain}

For $\PO$, \emph{neither} direction is forced:
$\Comp(\PP,\PO)=\{\PP,\PO,\DR\}$ says earlier positions are free;
$\Comp(\PPI,\PO)=\{\PPI,\PO\}$ says later ones are. A $\PO$-witness at
one phase need not be a $\PO$-witness at any other. And this is not a
hypothetical weakness of the proof --- the two-tier paper exhibits a
satisfiable concept that \emph{provably defeats} the constant-interface
architecture.

Take a chain with alternating types, $\tau_A$ at even positions
demanding $\exists\PO.A$, $\tau_B$ at odd positions demanding
$\forall\PO.\neg A$. Satisfiable? Yes --- by \textbf{wandering
witnesses}. For each $k$, witness $w_k$ (satisfying $A$) is:
\[
\begin{aligned}
\DR  &\ \text{ to the chain positions below } 2k
      &&\text{(disjoint from the small ones)}\\
\PO  &\ \text{ at exactly position } 2k
      &&\text{(overlaps that one)}\\
\PPI &\ \text{ to the positions above } 2k
      &&\text{(swallowed by the big ones)}
\end{aligned}
\]
Each $w_k$ slides through the chain: disjoint, then transiently
overlapping \emph{exactly one} rung, then swallowed forever. Every even
rung gets its $\exists\PO.A$ witness; every odd rung has \emph{no}
$\PO$-neighbor at all, so $\forall\PO.\neg A$ holds vacuously; all the
transitions are exactly the monotone ones from the table above; and the
witnesses are pairwise $\DR$ by forced composition. A perfectly good
infinite model.

But no \emph{constant-interface} certificate describes it: a witness
with recorded edge ``kernel${}\leftrightarrow w:\PO$'' would be $\PO$
to the $\tau_B$-rungs too, violating $\forall\PO.\neg A$. The demand at
every phase is real, but no \emph{single} witness serves two phases ---
the model needs infinitely many transient witnesses, and the finite
summary has no slot for ``transient.'' The architecture is provably
incomplete here. (Contrast $\DR$: the analogous villain with
$\exists\DR.A\,/\,\forall\DR.\neg A$ is \emph{unrealizable} ---
backward forcing drags any late $\DR$-witness back across the
$\forall\DR.\neg A$ rungs and clashes. The gap is a $\PO$-only
phenomenon; the table itself closes it for $\DR$.)

\subsection{Why $\forall\PO$-freeness dissolves the
gap}\label{ssec:dissolve}

Look at what actually went wrong: the constant-$\PO$ edge was fine
\emph{algebraically} (a constant-$\PO$ interface closes every triangle
along the chain --- check: $\PO\in\Comp(\PP,\PO)$ and
$\PO\in\Comp(\PPI,\PO)$, so the constant value is always among the
permitted ones) and failed only \emph{logically} --- it collided with a
$\forall\PO$ obligation at another phase. The Safe test for a
$\PO$-edge between types $\tau$ and $\sigma$ has exactly two parts: the
$\forall\PO$-universal part (``every $\forall\PO.E\in\tau$ must have
$E\in\sigma$,'' and symmetrically), and propositional bookkeeping.

\textbf{In a $\forall\PO$-free concept, no type contains any
$\forall\PO.E$.} The $\forall\PO$-universal part of every Safe test is
vacuous --- at every phase, for every witness. A witness type
containing $D$ is $\PO$-safe for \emph{all} phases the moment it is
$\PO$-safe for one (``automatic PO-coherence,'' Lemma~2 of the fragment
note). The one relation the table refuses to transport is also the one
relation nobody is watching. Constant-$\PO$ edges are algebraically
consistent and logically unchallenged; the wandering villain needs a
$\forall\PO$ to exist, and there are none.

\subsection{Counting, and the finish line}\label{ssec:count}

\begin{itemize}
\item Descriptors: $\le 2^{2^m}$. Witnesses: at most one per demand per
  descriptor, $\le(1+m)$ each. Quotient size
  $Q\le(1+m)\,2^{2^m}$; live width
  $K(C_0)\le Q\,(1+m)=(1+m)^2\,2^{2^m}$. Computable. (Nobody claims it
  is tight --- only computability matters, and no complexity bound is
  asserted.)
\item \textbf{Completeness:} extract from any model --- record
  recurrent types along each chain as descriptors, pick witnesses past
  the exhaustion index, let stabilization $+$ backward forcing $+$
  forward absorption $+$ automatic PO-coherence make them all-phase.
  Everything repeated folds into the kernels.
\item \textbf{Soundness:} unfold each kernel into its infinite chain
  under the constant interface, and prove the unfolding is still
  composition-closed on \emph{every} triple --- including triples whose
  closure jointly forces a $\PO$ edge (the multi-path caveat of
  \S\ref{sec:algebra} is enforced, not assumed away). This
  \textbf{chain-unfolding lift} is the soundness crux, and it is the
  part that is now \textbf{kernel-checked}: \code{POFreeLift.lean}
  proves, in Lean~4 with zero sorries, that if the finite augmented
  network is composition-consistent and strong-$\EQ$, then replacing
  the kernel by an infinite $\PP$-chain under the constant interface
  preserves composition-consistency. And under abstract semantics a
  composition-closed total network \emph{is} a model
  (\S\ref{sec:architecture}) --- so a valid certificate delivers an
  actual infinite model outright.
\end{itemize}

Enumerate certificates up to $K(C_0)$, check, done.

\section{Route 2: the ordered-disjoint normal form (the structural
argument)}\label{sec:route2}

{\small(The regular-cover pivot, GPT-5.5, July 2026:
\code{papers/rcc5\_local\_amalgamation\_theory.tex}, probes
\code{wp47}--\code{wp49}, \code{wp82}/\code{wp83}; normal form
certified in \code{formal/RCC5NormalForm.lean}, both directions.)}

Route 1 fought the composition table cell by cell. Route 2 wins by
changing the data structure so the fight never starts.

\subsection{A model is just an order and a disjointness
relation}\label{ssec:nf}

\textbf{Certified normal form.} A strong-$\EQ$ atomic RCC5 network is
composition-closed \textbf{iff} it is an \emph{ordered-disjoint
structure}:
\begin{itemize}
\item ${<}$ (that is, $\PP$) is a strict partial order;
\item ${\#}$ (that is, $\DR$) is symmetric, irreflexive, and
  \textbf{downward closed}: if $x\#y$, $x'\le x$, $y'\le y$, then
  $x'\#y'$;
\item comparable pairs are never disjoint;
\item \textbf{$\PO$ is the residual} --- it is \emph{defined} as
  ``neither comparable nor disjoint.'' It is not data. It is what's
  left.
\end{itemize}
(Forward direction kernel-checked for arbitrary domains in
\code{RCC5NormalForm.lean}; converse certified via an explicit
canonical set representation, cross-checked exhaustively to size four
and beyond, \code{wp47}/\code{wp88}.)

Let that reframing land: sixteen composition-table cells collapse into
two familiar closure laws --- transitivity of an order, downward
closure of a disjointness --- plus one compatibility condition. To
build a model you never ``place'' a $\PO$ edge, ever. You build an
order, you build a disjointness relation, you check two closure laws,
and every pair you said nothing about \emph{is} $\PO$, automatically,
consistently. The wild relation of \S\ref{sec:algebra} turns out to be
bookkeeping for ``no constraint here.''

With it come two machine-checked amalgamation facts we'll need:
\begin{itemize}
\item \textbf{Free amalgamation} (\code{wp48}): two ordered-disjoint
  structures glued along a shared part, with no new order or
  disjointness across, form an ordered-disjoint structure. Cross pairs
  land on the residual --- $\PO$.
\item \textbf{All-cross $\PO$ is always safe} (\code{wp82}): declaring
  \emph{every} cross pair of two closed pockets $\PO$ is always
  composition-closed.
\end{itemize}

\subsection{The tableau, and why it terminates}\label{ssec:tableau}

Now think like a tableau designer. Your nodes will form the vertical
skeleton (a tree of ${<}$-edges, one tree per component --- the
split-forest discipline); your constraints are $({<},{\#})$; $\PO$ is
whatever remains. Which obligations can arise, and how do they move?

Universals over $\PP,\PPI,\DR$ are exactly constraints on ${<}$ and
${\#}$: ``everything below me satisfies $D$,'' ``everything above me
satisfies $D$,'' ``everything disjoint from me satisfies $D$.'' How do
they propagate when edges compose? The general rule: if $x$ carries
$\forall R.D$ and $x\xrightarrow{\,S\,}y$, then $y$ must carry
$\forall T.D$ whenever $\Comp(S,T)\subseteq\{R\}$ (any $T$-successor of
$y$ is then a forced $R$-successor of $x$). Instantiating against the
table yields the \emph{complete} rule set (probe \code{wp49},
machine-checked from set semantics):

\begin{center}
\begin{tabular}{@{}lll@{}}
$\forall\PP.D$  & across $\PP$:  & require $D$, $\forall\PP.D$
  \quad(transitivity of ${<}$)\\
$\forall\PPI.D$ & across $\PPI$: & require $D$, $\forall\PPI.D$
  \quad(transitivity, downward)\\
$\forall\DR.D$  & across $\PP$:  & require $\forall\DR.D$
  \quad(my container's disjointnesses rain down onto me ---\\
  & & \phantom{require }so my container inherits my duty)\\
$\forall\DR.D$  & across $\DR$:  & require $D$, $\forall\PPI.D$
  \quad(my $\DR$-neighbor's parts are also my $\DR$-neighbors)
\end{tabular}
\end{center}

--- four rules, all riding the four singleton cells, and, decisively:

\begin{quote}
\textbf{No propagation rule from a $\forall\PP$, $\forall\PPI$, or
$\forall\DR$ obligation ever produces a $\forall\PO$ obligation.}
(Checked: no $\Comp(S,\PO)$ is contained in $\{\PP\}$, $\{\PPI\}$, or
$\{\DR\}$.)
\end{quote}

So the obligation vocabulary
$\{\forall\PP,\forall\PPI,\forall\DR\}\times\cl(C_0)$ is \textbf{closed
under propagation}. Start $\forall\PO$-free, stay $\forall\PO$-free ---
not just syntactically in the input, but dynamically, in every label
the tableau will ever write. Obligations are drawn from a fixed finite
set, and they move \emph{deterministically} along the singleton cells
(the same four cells as Route~1, wearing different clothes:
transitivity of ${<}$, downward closure of ${\#}$). Node labels $=$
type $+$ obligation set $=$ finitely many possibilities
$\Rightarrow$ standard subformula blocking on the vertical tree
terminates, with the loops read as generators
(\S\ref{sec:architecture}).

What about the demands?
\begin{itemize}
\item \textbf{$\exists\PP\,/\,\exists\PPI\,/\,\exists\DR$:} create the
  witness in the skeleton or as a sibling; its forced consequences are
  computed by the two closure laws --- finitely, deterministically
  (down-closure of ${\#}$ is a forced flood, but a flood of
  \emph{shadows}: recoverable, costless, per \S\ref{sec:architecture}).
\item \textbf{$\exists\PO.D$:} here is the beautiful step. Build a
  satisfying structure for $D$ separately (recursion), then glue it on
  declaring \textbf{no cross order and no cross disjointness at all}
  --- free amalgamation. Every cross pair, \emph{including the
  demanding node and its witness}, lands on the residual: $\PO$. The
  demand is satisfied by the default itself. Is that legal?
  \emph{Algebraically} always (all-cross $\PO$ is always safe,
  \code{wp82}). \emph{Logically} --- a cross-$\PO$ edge could only
  violate a $\forall\PO$ obligation, \textbf{and there are none, and
  there never will be any} (closure of the vocabulary). Both halves of
  ``safe'' hold, and the second half is exactly $\forall\PO$-freeness.
\end{itemize}

That is the ``yes!'' moment for Route~2: \textbf{the always-consistent
default policy (make everything unconstrained overlap) is also always
obligation-free --- precisely in this fragment.} The full logic's
nightmare --- steering $\PO$ edges jointly against watching universals
--- never engages, because nothing ever watches. The horizontal axis
carries no obligations; all real work happens on the vertical skeleton,
where the singleton cells make propagation deterministic and blocking
sound.

Status, honestly: Route~2 is a \emph{re-derivation} --- its ingredients
(the normal form, the propagation table, the amalgamation and
cross-policy facts) are individually machine-checked, two of them
kernel-certified, but the assembly into a standalone
tableau-completeness proof is theorem-level prose. It independently
confirms the fragment result; it does not independently certify it.

\section{The boundary is exactly right --- and it is not ``horizontal
vs.\ vertical''}\label{sec:boundary}

A skeptical reader should now ask: is $\forall\PO$-freeness really the
boundary, or just a convenient place to stop? Four checks.

\paragraph{The fragment keeps a horizontal universal.} Note what did
\emph{not} have to go: $\forall\DR$. Disjointness is as ``horizontal''
as overlap, yet $\forall\DR.D$ is harmless. The dividing line is not
horizontal-vs-vertical but \textbf{stable vs.\ unstable under vertical
motion}: $\DR$ is transported deterministically by the singleton cells
(rains down into parts); $\PO$ is transported by no singleton cell in
either direction and is never the forced outcome of any single
composition. The fragment excises the unstable relation's universal and
nothing else. That the cut lands on a single constructor is a fact
about the composition table, not a choice.

\paragraph{The gap on the other side is real.} The wandering-witness
concept of \S\ref{ssec:villain} is a \emph{satisfiable} concept just
outside the fragment (one $\forall\PO$) that provably defeats the
constant-interface certificate; the probe \code{wp86} (Part~C) confirms
the boundary computationally --- the PO-incoherent descriptor forces an
unsatisfiable interface, the coherent ones never do. And the full
logic's signature satisfiable diagnostic, the strong-$\EQ$ $\PO$-loop
$C\sqcap\exists\PO.\exists\PO.C\sqcap\forall\{\PO,\DR,\PP,\PPI\}.\neg C$,
needs its $\forall\PO$ conjunct to force the loop closure --- it, too,
lives outside. Everything that has ever been \emph{hard} here has a
$\forall\PO$ in it.

\paragraph{The connection to the open problem is exact.} The full-logic
keystone F6 asks for a computable live-width bound. Both routes above
\emph{prove} one for the fragment (Route~1 explicitly: $K(C_0)$). And
the fragment theorem is itself the sharpest available evidence about
where the full problem's difficulty lives: remove \emph{one}
constructor --- the universal over the one relation the singleton cells
don't control --- and the coordination problem that every full-logic
attempt has broken on dissolves. ($\forall\DR$ also coordinates
horizontal edges, but its coordination is \emph{deterministic}, which
is exactly why it may stay.) The fragment is thus not a lucky island;
it is the complement of the obstruction, and its proof shows \emph{why}
the obstruction is where it is.

\paragraph{A corollary: in the fragment, identity can never be forced.}
The project's knife's-edge analysis (the coincidence obstruction of the
overview paper) turns on \emph{merge-forcing}: clash out all
non-$\EQ$ options between two witnesses, and $\EQ$ --- identity, under
strong $\EQ$ --- is all that remains. Such forced merges are the raw
material of grid corners and rigid addressing, the machinery every
undecidability attempt needs. In the fragment the trick is
inexpressible, by the same escape-valve fact one level up:
machine-checked (\code{wp14}, part~T4), \textbf{every composition cell
containing $\EQ$ also contains $\PO$} --- and hence so does every
\emph{intersection} of such cells (if $\EQ$ survives a joint
constraint, every constituent contained $\EQ$, so every constituent
contained $\PO$, so $\PO$ survives too). Wherever identity is on a
pair's menu, partial overlap is on it as well, and the fragment has no
eraser for $\PO$. No $\forall\PO$-free concept ever forces two nodes to
be one. So the fragment sits provably on the ``loose'' side of the
knife's edge: not only does the decidability proof close --- the
undecidability \emph{machinery} cannot even be assembled.

\section{So --- was it proved twice, or three times?}\label{sec:ledger}

Honest ledger. There are \textbf{two genuinely different argument
shapes} and \textbf{three arrivals}, plus one certification pass.

\begin{enumerate}
\item \textbf{April 2026 --- the two-tier quotient} (Claude, reviewed
  through four revisions by GPT-5.4). Proves decidability for the
  \emph{PO-coherent} fragment --- concepts whose $\forall\PO$
  obligations, if any, are uniform across phases --- with
  $\forall\PO$-free concepts as the clean special case, and honestly
  exhibits the $\PO$ gap (\S\ref{ssec:villain}) as the boundary.
  Route~1 above.
\item \textbf{July 2026 --- the cold F6/W2$'$ attack} (GPT-5.5, given
  the open problem cold, asked to attack it). Independently arrived at
  the $\forall\PO$-free fragment as its ``Target B'' deliverable, with
  the $K(C_0)$ bound (\code{papers/po\_free\_fragment\_ALCIRCC5.tex}).
  On inspection its algebra --- stabilization, backward forcing,
  forward absorption, automatic PO-coherence --- is the two-tier
  argument's, so it counts as an independent \emph{re-derivation} of
  Route~1, not a new route: the same theorem found again by a different
  mind from a different starting point, which is worth something
  evidentially but is not a second proof shape.
\item \textbf{July 2026 --- the regular-cover pivot} (GPT-5.5, later
  phase). Re-derives the fragment through the ordered-disjoint normal
  form and the propagation-closure fact (\code{wp49}) --- Route~2
  above, a genuinely different shape: no chains, no phases, no
  stabilization; instead a change of data structure that makes $\PO$
  disappear as data.
\item \textbf{2026-07-18 --- the certification pass.} After a cold
  review flagged that the fragment's original rationale leaned on an
  overstated ``$\PO$ is never forced'' (the calibration of
  \S\ref{sec:algebra}), the fragment was re-examined: it holds, the
  argument enforces full composition closure, and the soundness crux
  --- the chain-unfolding lift --- was kernel-checked
  (\code{POFreeLift.lean}, zero sorries). Additionally, two
  \emph{independent} decision procedures (the cover-tree tableau and
  the quasimodel reasoner) were run on 244 random and curated
  $\forall\PO$-free concepts with zero mismatches (\code{wp87}) --- a
  clean cross-validation, because the quasimodel oracle's one known
  blind spot needs a $\forall\PO$ and so cannot arise in the fragment.
\end{enumerate}

\paragraph{A gap in Route~1, found 2026-08-28 --- and what it does not
touch.} Re-reading the two-tier paper turned up a genuine defect in its
completeness argument: Step~2 defined the relation between two kernels as a
double limit $\lim_{i,j\to\infty}\rho(d_i,d_j)$, justified by external-relation
stabilization ``applied twice''. That lemma is \emph{one-sided}---it fixes an
element and varies the chain---and the double limit need not exist
($\mathbb{Z}\times\{0,1\}$ ordered by first coordinate: every row and every
column stabilizes, yet every tail carries $\PP$, $\PO$ and $\PPI$). It is
repaired by a \emph{simultaneous finite rectangle} in place of the limit, which
is all the quotient needs since it reads only finitely many phases; the repair is
machine-checked (\code{fused\_kq\_all}).

Which arrivals this touches is worth stating precisely, because it is the
argument for keeping two routes:
\begin{itemize}
\item \textbf{Arrival~1 (Route~1)} stated the false step. Repaired.
\item \textbf{Arrival~2} re-derives Route~1 and uses external-relation
  stabilization \emph{one-sidedly}, as the lemma allows; it never states the
  double limit. It builds a two-tier quotient, so it depends on the repaired
  lemma for its kernel-to-kernel edges, but it did not assert the invalid
  inference.
\item \textbf{Arrival~3 (Route~2)} is \emph{untouched}. It has no chains, no
  phases and no stabilization at all---the defect cannot be stated in its
  vocabulary, let alone used.
\end{itemize}
So the theorem never rested on the defective step: an independent route to it
never had the defect. This is exactly the redundancy two proof shapes are for.

So: ``2 (or 3?)'' resolves to \textbf{two independent proof routes,
three independent arrivals at the theorem} --- and the fact that a cold
attacker, pointed at the open problem with no steer, walked to this
exact fragment on its own is perhaps the best evidence that the
boundary is natural.

\section{What is certified, what is checked, what is
argued}\label{sec:status}

\begin{center}
\small
\begin{tabular}{@{}>{\raggedright\arraybackslash}p{0.52\textwidth}>{\raggedright\arraybackslash}p{0.42\textwidth}@{}}
\toprule
\textbf{Claim} & \textbf{Status}\\
\midrule
Composition/converse tables; the singleton and $\PO$ cells of
\S\ref{sec:algebra} & Kernel-certified (Lean, \code{decide}) and
derived from set semantics in every probe\\[0.3em]
RCC5 normal form (network $\Leftrightarrow$ ordered-disjoint),
arbitrary domains & \textbf{Certified}, both directions
(\code{RCC5NormalForm.lean} $+$ canonical representation)\\[0.3em]
Chain-unfolding lift (finite CC certificate $\Rightarrow$ CC infinite
unfolding) & \textbf{Certified} (\code{POFreeLift.lean}, zero
sorries)\\[0.3em]
Certificate-to-model soundness: a valid single-kernel two-tier
certificate yields a model of the concept, through the logic (Hintikka
$+$ truth lemma on the unfolding) & \textbf{Certified}
(\code{POFreeLift.lean} round~A, 2026-07-22: \code{twoTier\_sound};
witness \code{cinf\_satisfiable} --- the no-finite-model concept
$\exists\PP.\top\sqcap\forall\PP.\exists\PP.\top$, satisfiable
end-to-end)\\[0.3em]
Multi-kernel certificates (any family of ascending/descending kernels)
and the executable first-order checker (certificates as pure list
data; acceptance $\Rightarrow$ model) & \textbf{Certified}
(\code{POFreeLift.lean} rounds~B--C: \code{multiTier\_sound},
\code{mtAcceptB\_sound}; the ascending$+$descending two-tower witness
accepted by kernel \code{decide})\\[0.3em]
Propagation closure (no non-$\PO$ universal creates a $\forall\PO$
obligation) & Machine-checked exhaustively (\code{wp49})\\[0.3em]
Free amalgamation; all-cross-$\PO$ safety; non-uniform cross-policies
& Machine-checked exhaustively (\code{wp48}, \code{wp82},
\code{wp83})\\[0.3em]
Lift-lemma stress-test incl.\ multi-path $\PO$ forcing; fragment
boundary & Machine-checked empirically (\code{wp86})\\[0.3em]
Two independent reasoners agree on the fragment (244 concepts) &
Machine-checked empirically (\code{wp87})\\[0.3em]
\textbf{Fragment decidability itself}, by a \emph{third} route (a finite
control graph of signatures, a monotone elimination to a greatest fixed
point, and a fresh-occurrence unfolding) --- on raw input and under
concrete set semantics & \textbf{Certified} (2026-08-29,
\code{POFreeLift.lean}: \code{decidableFSat}, \code{decidableSat\_cone},
\code{decidableSetSat}; fragment membership the only hypothesis, zero
sorries). Cold-reviewed three times without a counterexample; \emph{not}
certified in this project's stricter sense --- see the ledger. Decidable,
\emph{not} runnable: the signature space is doubly exponential\\[0.3em]
The two \emph{arguments} of this paper ($K(C_0)$ bound, extraction,
assembly --- Routes 1 and 2) & \textbf{Theorem-level} --- argued, twice,
independently; \emph{not} end-to-end formalized. The certified route above
is a \emph{different} proof of the same theorem and does not use
$K(C_0)$\\
\bottomrule
\end{tabular}
\end{center}

Remaining for full certification --- the campaign is underway: rounds
A--D1 (2026-07-22/23) landed the entire \textbf{soundness side plus
the decision architecture}: a valid two-tier certificate, with any
family of ascending or descending kernels, presented as \textbf{pure
first-order list data}, is accepted by a computable oracle-free
Boolean checker \textbf{exactly when valid} (\code{mtOkB\_iff}),
acceptance yields a model (kernel-checked end to end; the two-tower
witness's acceptance runs inside the kernel via plain \code{decide}),
and a certified reduction (\code{decidableSat\_of\_codes}) turns any
candidate-code list plus a completeness premise into
$\mathrm{Decidable}$ satisfiability. Rounds D2a--c further certified
the extraction's entire \textbf{model-side toolkit}: the stabilization
theorem of \S\ref{ssec:crux} with both forcing corollaries, the
infinite pigeonhole (recurrent tails, segment selection), segment
coherence (cycling a type-equal chain segment into a kernel is
legitimate), witness selection (late picking), and the syntactic
$\forall\PO$-vacuity (\code{pofree\_cl\_all}: no $\forall\PO$ ever
appears in the fragment's closure --- the escape valve,
kernel-checked). Still open: the assembly construction gluing these
into an accepted code, and the $K(C_0)$ counting/enumeration (see
\code{LEAN.md} for the recorded design). Until then, the honest label is the project's
standing one: \textbf{strongly supported, with the soundness core
machine-certified --- not certified end-to-end.}

\section{The one-paragraph version}\label{sec:summary}

$\ALCIRCC$ models are complete graphs with a composition law, and they
are legitimately infinite, so deciding satisfiability means enumerating
and checking \emph{finite certificates} --- blueprints whose loops are
generators of infinite models --- and the whole difficulty is bounding
the certificate's genuinely-free (``live'') content. The composition
table transports $\DR,\PP,\PPI$ deterministically along the nesting
order (four singleton cells: transitivity, and disjointness raining
down into parts) but refuses to transport $\PO$ in either direction,
and no single composition ever forces $\PO$. Consequently universal
obligations over $\DR/\PP/\PPI$ are trackable --- they propagate
deterministically in a closed finite vocabulary --- while $\forall\PO$
obligations must chase neighbors the algebra refuses to pin down; a
concrete satisfiable concept with one $\forall\PO$ provably defeats the
finite constant-interface certificate via ``wandering witnesses'' that
each overlap a chain only transiently. Ban $\forall\PO$ and both known
proof shapes close: the \emph{two-tier quotient} folds each infinite
chain into a kernel whose witnesses' relations stabilize and are
all-phase by forced composition (with constant-$\PO$ edges
algebraically consistent and, now, unchallenged), giving a computable
live-width bound and enumerate-and-check decidability; the
\emph{ordered-disjoint normal form} --- certified --- shows a model is
nothing but a strict order plus a downward-closed disjointness with
$\PO$ as the residual non-constraint, so a standard tree tableau
tracking only order/disjointness obligations (whose propagation
calculus provably never manufactures a $\forall\PO$ obligation)
terminates by ordinary blocking, discharging $\exists\PO$ by
always-safe free amalgamation. Two shapes, three independent arrivals,
the soundness crux kernel-checked in Lean; the fragment is exactly the
region where the open keystone F6 --- bound the live width --- is a
theorem instead of a conjecture, and its boundary, one constructor
wide, is drawn by the composition table itself.

\bigskip
\begin{center}
\fbox{\parbox{0.92\textwidth}{\small\itshape
Companions: \code{WHY\_ITS\_HARD.md} (the full problem, non-technical),
\code{papers/two\_tier\_quotient\_ALCIRCC5.tex} (Route~1 in full),
\code{papers/po\_free\_fragment\_ALCIRCC5.tex} (the fragment note),
\code{papers/rcc5\_local\_amalgamation\_theory.tex} (Route~2's local
algebra), \code{papers/overview\_arxiv.tex}, \S``The win'' (the
calibrated summary), \code{formal/POFreeLift.lean},
\code{formal/RCC5NormalForm.lean} (the certified parts). The Markdown
original of this note is \code{papers/WHY\_PO\_FREE\_IS\_DECIDABLE.md}.
Where this document and the formal artifacts disagree, the formal
artifacts win.}}
\end{center}

\end{document}
